\section*{3. 孤立奇点}

\subsection*{3.1. 消灭圈流形}

\subsubsection*{3.1.1.}
设
\[
D=\{z\in\mathbb C\mid |z|<1\},
\]
设 $X$ 是解析空间，$0\in X$，且 $f:X\to D$ 是满足 $f(0)=0$ 的态射。假设 $f$ 在 $0$ 之外光滑（特别地，$X$ 在 $0$ 之外光滑）。我们要研究 $0$ 附近的消灭圈；因此，不失一般性，可假设 $X$ 嵌入 $\mathbb C^n\times D\subset\mathbb C^{n+1}$，与到 $D$ 的投影相容，并将 $0$ 映至 $0$。

\subsubsection*{3.1.2.}
设 $\delta(x)$ 是 $X$ 上的实解析函数（即 $\mathbb C^{n+1}$ 上实解析函数在 $X$ 上的限制），满足对 $x\in X_0=f^{-1}(0)$，

\begin{enumerate}[label=(\alph*)]
\item $\delta(x)\geq 0$；
\item $\delta(x)=0$ 当且仅当 $x=0$。
\end{enumerate}

例如，可取 $\delta$ 为 $\mathbb C^{n+1}$ 中 $x$ 到 $0$ 的 Euclid 距离的平方。$\delta$ 在 $X_0$ 上的限制的临界点集是实解析的，因此在 $0$ 的适当邻域内只有有限个连通分支。将 $X$ 换为 $0$ 的开邻域，并应用 Sard 定理，于是可达到如下情形：

\begin{enumerate}[label=(\alph*),start=3]
\item 对适当的 $r>0$，空间 $X_0\cap\delta^{-1}([0,r])$ 紧；
\item 在 $X_0$ 上，$\delta$ 在 $(0,r]$ 中没有临界值。
\end{enumerate}

进一步缩小 $X$ 后，还有：

\begin{enumerate}[label=(\alph*),start=5]
\item 存在一个 $C^\infty$ 函数 $\tau(s)$，$0\leq s\leq r$，满足 $\tau(0)=0$ 且对 $s>0$ 有 $\tau'(s)>0$，使得对 $0<s\leq r$ 及 $0\leq |t|^2\leq\tau(s)$，纤维 $X_t=f^{-1}(t)$ 与 $X$ 中的实超曲面 $\delta(x)=s$ 横截，且盒形区域
\[
  B_{t,d}=\{x\in X\mid \delta(x)\leq d\ \text{且}\ |f(x)|^2<t\}
  \]
在半径为 $t^{1/2}$ 的开圆盘 $D_t$ 上固有。
\end{enumerate}

设 $L=X_0\cap\delta^{-1}(r)$。由 (c) 与 (d)，存在从 $X_0\cap\delta^{-1}([0,r])$ 到锥体
\[
L\times[0,r]\big/\bigl(L\times\{0\}\text{ 压缩为一点}\bigr)
\]
的同胚，将 $\delta$ 化为第二个坐标，且在原点外为 $C^\infty$ 的。

\subsubsection*{3.1.3.}
设 $d\leq r$ 及 $t\leq\tau(d)$。在 $0$ 之外，$B_{t,d}$ 是带角流形，投影 $f:B_{t,d}\to D_t$ 在 $D_t^*$ 上为纤维化。\EditorialFootnote{原文将此投影直接称为纤维化。这里明确限定到穿孔圆盘：随后 (a) 给出的正是在 $D_t^*$ 上的光滑丛；(b) 则说明其边界丛可以延拓到 $D_t$。一般不能把奇异纤维在 $0$ 处的局部平凡性也包括在内。}特别地：

\begin{enumerate}[label=(\alph*)]
\item 在 $D_t^*=D_t-\{0\}$ 上，$B_{t,d}$ 诱导一个紧带边界流形的 $C^\infty$ 丛；
\item 此丛的边界 $\delta(x)=d$ 延拓为 $D_t$ 上紧流形的 $C^\infty$ 丛 $\partial B_{t,d}$。
\end{enumerate}

设 $y\in D_t^*$。称带边界流形 $V=f^{-1}(y)$ 为\underline{消灭圈流形}。由于圆盘可缩，$D_t$ 上的丛 $\partial B_{t,d}$ 同构于以 $\partial V$ 为纤维的常值丛，并且在 $y$ 上的纤维上限制为恒等映射的丛同构
\[
\alpha:\partial B_{t,d}\xrightarrow{\ \sim\ }\partial V\times D_t
\]
所构成的空间可缩。选择其中一个，记为 $\alpha$。将半径为 $|y|$ 的圆周 $S_{|y|}$ 上的向量场 $\partial/\partial\theta$ 提升为 $f^{-1}(S_{|y|})$ 上的向量场 $u$，使其在边界上与 $\alpha$ 相容，即
\[
d\alpha(u)=(0,\partial/\partial\theta).
\]
这样的提升所成的空间可缩。将 $u$ 沿正向绕圆周积分一周，定义 $V$ 到自身的微分同胚 $T$，它在 $\partial V$ 上为恒等；这就是\underline{单值变换}。它在相对于 $\partial V$ 的同痕意义下良定义。

在一个留待读者精确表述的意义下，$(V,T)$ 不依赖于 $t$ 和 $d$ 的选择。我不知道它是否不依赖于 $\delta$ 的选择。不过，若要求 $\delta$ 是 $\mathbb C^{n+1}$ 上实解析函数 $g$ 在 $X$ 上的限制，并且在 $X_0$ 于 $0$ 处的 Zariski 切空间上满足 $dg=0$ 及 $d^2g>0$（正定），则确实不依赖于这一选择。事实上，对这样的 $\delta$ 和 $\delta'$，可像 3.1.2(c),(d) 中一样选择 $r$，使之对每个 $\lambda\delta+(1-\lambda)\delta'$，$0\leq\lambda\leq1$ 一致成立。

